\documentclass[11pt]{article}
\usepackage{fontspec,amsmath,amssymb,amsthm,hyperref}
\usepackage[margin=1in]{geometry}
\setlength{\emergencystretch}{3em}
\title{Integrate a forced entropy-dissipation inequality}
\author{LangevinLean mathematical exposition}
\date{}
\begin{document}\maketitle
\noindent Source module: \texttt{UnadjustedLangevin/LSI/EntropyDecay.lean}. Compile this proof document with XeLaTeX or LuaLaTeX. The accompanying HTML pairs the mathematical argument with the exact Lean statements and records the scope of the formalization.


An integrating factor solves \(H\prime\le-aH+b\) as \(H(t)\le e^{-at}H(0)+(b/a)(1-e^{-at})\), under the displayed positive-rate, continuity and one-sided derivative assumptions. The entropy wrapper supplies the appropriate rate from LSI and dissipation.


\section{A scalar affine comparison from right derivatives}


For continuous \(H\) with right derivative bounded by \(-aH+b\), \(a>0\), compare it to the exact scalar ODE solution. The one-sided Grönwall theorem accepts the right derivative through its right-slope liminf inequality. Simplifying its explicit solution gives \(H(t)\le e^{-at}H(0)+(b/a)(1-e^{-at})\). Neither \(H\) nor \(b\) needs to be nonnegative.



\noindent Lean declarations: \texttt{affine\_decay\_of\_deriv\_right\_le}.

\section{Scalar LSI and dissipation give rate 2\ensuremath{\alpha}}


The supplied dissipation identity is \(H'=-I\), and the scalar LSI premise is \(2\alpha H\le I\). Therefore \(H'\le-2\alpha H\). Apply the affine comparison theorem with \(a=2\alpha\) and \(b=0\), obtaining \(H(t)\le e^{-2\alpha t}H(0)\). This scalar theorem does not itself establish the dissipation identity for a law-valued evolution.



\noindent Lean declarations: \texttt{entropy\_decay\_of\_dissipation}.

\end{document}
